\documentclass[12pt,a4paper]{article}

\usepackage[T2A]{fontenc}
\usepackage[utf8]{inputenc}
\usepackage[russian]{babel}
\usepackage{amsmath}
\usepackage{amssymb}
\usepackage{geometry}
\usepackage{listings}
\geometry{margin=2.5cm}

\title{Факторизованная грамматика для продвинутого лямбда-интерпретатора}
\author{}
\date{}

\begin{document}

\maketitle

\section*{Условие}
Требуется построить факторизованную однозначную грамматику без
$\varepsilon$-переходов для языка лямбда-выражений с именованными
определениями вида \texttt{let}. Поддерживаются:
\begin{itemize}
    \item переменные;
    \item аппликация;
    \item лямбда-абстракция;
    \item скобки;
    \item несколько параметров у лямбда-абстракции;
    \item набор именованных определений перед основным выражением.
\end{itemize}

\section*{Терминалы}
Будем считать, что лексический анализатор выделяет следующие токены:
\begin{itemize}
    \item \texttt{LET} --- ключевое слово \texttt{let};
    \item \texttt{LAM} --- символ \texttt{\textbackslash};
    \item \texttt{DOT} --- символ \texttt{.};
    \item \texttt{EQ} --- символ \texttt{=};
    \item \texttt{LPAREN}, \texttt{RPAREN} --- скобки \texttt{(} и \texttt{)};
    \item \texttt{IDENT} --- идентификатор (кроме ключевого слова \texttt{let});
    \item \texttt{NL} --- конец строки.
\end{itemize}

\section*{Грамматика}

\begin{align*}
\textit{Program} &::= \textit{Definitions}\ \textit{Expr}
                  \mid \textit{Expr} \\
\textit{Definitions} &::= \textit{Definition}
                      \mid \textit{Definition}\ \textit{Definitions} \\
\textit{Definition} &::= \texttt{LET}\ \texttt{IDENT}\ \texttt{EQ}\ \textit{Expr}\ \texttt{NL} \\
\textit{Expr} &::= \textit{Lambda}
               \mid \textit{Application} \\
\textit{Lambda} &::= \texttt{LAM}\ \textit{Params}\ \texttt{DOT}\ \textit{Expr} \\
\textit{Params} &::= \texttt{IDENT}
                 \mid \texttt{IDENT}\ \textit{Params} \\
\textit{Application} &::= \textit{Atom}
                      \mid \textit{Atom}\ \textit{Application} \\
\textit{Atom} &::= \texttt{IDENT}
               \mid \texttt{LPAREN}\ \textit{Expr}\ \texttt{RPAREN}
\end{align*}

\section*{Свойства грамматики}

\begin{itemize}
    \item В грамматике отсутствуют $\varepsilon$-продукции.
    \item Грамматика не содержит левой рекурсии.
    \item Лямбда-абстракция однозначно распознаётся по начальному символу \texttt{\textbackslash}.
    \item Аппликация задаётся как непустая последовательность атомов.
\end{itemize}

Последовательность атомов
\[
a_1\ a_2\ \dots\ a_n
\]
интерпретируется как левоассоциативная аппликация:
\[
(((a_1\ a_2)\ a_3)\dots a_n).
\]

Лямбда-абстракция с несколькими параметрами
\[
\texttt{\textbackslash x y z.M}
\]
интерпретируется как
\[
\texttt{\textbackslash x.\textbackslash y.\textbackslash z.M}.
\]

\section*{Пример вывода 1}

Выведем выражение:
\[
\texttt{\textbackslash x y.x}
\]

\begin{align*}
\textit{Expr}
&\Rightarrow \textit{Lambda} \\
&\Rightarrow \texttt{LAM}\ \textit{Params}\ \texttt{DOT}\ \textit{Expr} \\
&\Rightarrow \texttt{LAM}\ \texttt{IDENT}\ \textit{Params}\ \texttt{DOT}\ \textit{Expr} \\
&\Rightarrow \texttt{LAM}\ \texttt{IDENT}\ \texttt{IDENT}\ \texttt{DOT}\ \textit{Expr} \\
&\Rightarrow \texttt{LAM}\ \texttt{IDENT}\ \texttt{IDENT}\ \texttt{DOT}\ \textit{Application} \\
&\Rightarrow \texttt{LAM}\ \texttt{IDENT}\ \texttt{IDENT}\ \texttt{DOT}\ \textit{Atom} \\
&\Rightarrow \texttt{LAM}\ \texttt{IDENT}\ \texttt{IDENT}\ \texttt{DOT}\ \texttt{IDENT}
\end{align*}

\section*{Пример вывода 2}

Выведем выражение:
\[
\texttt{S K K}
\]

\begin{align*}
\textit{Expr}
&\Rightarrow \textit{Application} \\
&\Rightarrow \textit{Atom}\ \textit{Application} \\
&\Rightarrow \texttt{IDENT}\ \textit{Application} \\
&\Rightarrow \texttt{IDENT}\ \textit{Atom}\ \textit{Application} \\
&\Rightarrow \texttt{IDENT}\ \texttt{IDENT}\ \textit{Application} \\
&\Rightarrow \texttt{IDENT}\ \texttt{IDENT}\ \textit{Atom} \\
&\Rightarrow \texttt{IDENT}\ \texttt{IDENT}\ \texttt{IDENT}
\end{align*}

\section*{Пример программы}

\begin{lstlisting}
let S = \x y z.x z (y z)
let K = \x y.x
S K K
\end{lstlisting}

\noindent
Первые две строки соответствуют нетерминалу \textit{Definitions}, последняя строка --- \textit{Expr}.

\end{document}